\section{Variational Autoencoder (VAE)}

\subsection{De la codificación determinista a la probabilística}

La codificación de un Autoencoder tradicional es determinista: dado un dato de entrada $x$, el codificador produce un vector latente fijo $z$. Esto puede hacer que el espacio latente tenga "huecos" y regiones discontinuas; muestrear desde esas regiones produce salidas sin sentido.

\begin{figure}[H]
\centering
\includegraphics[width=0.85\textwidth]{slides-images/slide-20.jpg}
\caption{Flujo de codificación determinista de un Autoencoder tradicional\protect\footnotemark}
\end{figure}
\footnotetext{Slides página 20.}

El \textbf{Variational Autoencoder (VAE)} introduce una mejora probabilística: el codificador ya no produce un vector latente fijo, sino que produce los \textbf{parámetros de una distribución de probabilidad}: un vector de medias $\mu$ y un vector de desviaciones estándar $\sigma$. Después se \textbf{muestrea} de esa distribución para obtener el vector latente $z$.

\begin{figure}[H]
\centering
\includegraphics[width=0.85\textwidth]{slides-images/slide-22.jpg}
\caption{La mejora central del VAE: espacio latente probabilístico\protect\footnotemark}
\end{figure}
\footnotetext{Slides página 22. Fuente: Kingma+ ICLR 2014.}

\begin{importantbox}{La idea central del VAE}
El VAE es una \textbf{mejora probabilística} del Autoencoder. El codificador produce una media $\mu$ y una desviación estándar $\sigma$, que definen una distribución de probabilidad sobre la variable latente $q_\phi(z|x)$. Al muestrear de esa distribución, el VAE puede generar distintos vectores latentes, logrando así una generación diversa.
\end{importantbox}

\subsection{La función de pérdida del VAE}

El objetivo de optimización del VAE se compone de dos partes:

\begin{figure}[H]
\centering
\includegraphics[width=0.85\textwidth]{slides-images/slide-25.jpg}
\caption{Optimización del VAE: pérdida de reconstrucción + término de regularización\protect\footnotemark}
\end{figure}
\footnotetext{Slides página 25. Fuente: Kingma+ ICLR 2014.}

$$\mathcal{L}(\phi, \theta, x) = \underbrace{\mathbb{E}_{q_\phi(z|x)}[\log p_\theta(x|z)]}_{\text{Pérdida de reconstrucción (Reconstruction Loss)}} - \underbrace{D_{KL}(q_\phi(z|x) \| p(z))}_{\text{Término de regularización (Regularization Term)}}$$

\begin{itemize}
    \item \textbf{Pérdida de reconstrucción}: mide la calidad con la que el decodificador reconstruye el dato $\hat{x}$ a partir de la variable latente $z$, y fomenta que $\hat{x}$ se acerque lo más posible a $x$
    \item \textbf{Término de regularización (KL Divergence)}: mide la distancia entre la distribución $q_\phi(z|x)$ de la variable latente producida por el codificador y la distribución previa $p(z)$, y fomenta que el espacio latente tenga una buena estructura
    \item $\phi$: los parámetros del codificador
    \item $\theta$: los parámetros del decodificador
\end{itemize}

\subsection{La distribución previa y la KL Divergence}

El VAE necesita fijar una \textbf{distribución previa} (Prior) $p(z)$ para la variable latente. La opción más común es la distribución normal estándar:

$$p(z) = \mathcal{N}(\mu = 0, \sigma^2 = 1)$$

\begin{figure}[H]
\centering
\includegraphics[width=0.85\textwidth]{slides-images/slide-28.jpg}
\caption{Distribución previa: la distribución normal fomenta que las variables latentes se distribuyan de manera uniforme alrededor del centro del espacio latente\protect\footnotemark}
\end{figure}
\footnotetext{Slides página 28. Fuente: Kingma+ ICLR 2014.}

\begin{knowledgebox}{¿Por qué elegir una distribución previa normal?}
Elegir $\mathcal{N}(0, 1)$ como distribución previa tiene dos propósitos importantes:
\begin{enumerate}
    \item \textbf{Codificación concentrada}: fomenta que el codificador distribuya las variables latentes en la región central del espacio latente
    \item \textbf{Prevenir "trampas"}: evita que la red agrupe puntos de datos de distintas clases en regiones aisladas específicas del espacio latente (es decir, que memorice los datos), y en su lugar la obliga a aprender una representación significativa y continua
\end{enumerate}
\end{knowledgebox}

La \textbf{KL Divergence} (divergencia de KL) es una medida de la "distancia" entre dos distribuciones de probabilidad. Para una distribución previa normal, la divergencia de KL tiene una expresión analítica:

$$D_{KL}(q_\phi(z|x) \| p(z)) = -\frac{1}{2} \sum_{j=1}^{J} \left(1 + \log(\sigma_j^2) - \mu_j^2 - \sigma_j^2\right)$$

donde $J$ es la dimensión de la variable latente, y $\mu_j$ y $\sigma_j$ son, respectivamente, la media y la desviación estándar de la $j$-ésima variable latente.

\subsection{La intuición de la regularización: continuidad e integridad}

La función de la regularización es garantizar que el espacio latente tenga dos propiedades clave:

\begin{figure}[H]
\centering
\includegraphics[width=0.85\textwidth]{slides-images/slide-30.jpg}
\caption{La regularización garantiza la continuidad y la integridad del espacio latente\protect\footnotemark}
\end{figure}
\footnotetext{Slides página 30. Fuente: Joseph Rocca.}

\begin{enumerate}
    \item \textbf{Continuidad (Continuity)}: puntos adyacentes en el espacio latente deberían producir, al decodificarse, contenidos semánticamente similares
    \item \textbf{Integridad (Completeness)}: muestrear desde cualquier posición del espacio latente debería decodificarse en contenido con sentido
\end{enumerate}

\begin{figure}[H]
\centering
\includegraphics[width=0.85\textwidth]{slides-images/slide-32.jpg}
\caption{Tras la regularización se garantizan la continuidad y la integridad del espacio latente\protect\footnotemark}
\end{figure}
\footnotetext{Slides página 32. Fuente: Joseph Rocca.}

\begin{warningbox}{El problema del espacio latente sin regularización}
Si no hay regularización (y solo se usa la pérdida de reconstrucción), el espacio latente que la red puede aprender tendrá numerosos "huecos", regiones que los datos de entrenamiento no cubren. Muestrear de las regiones vacías producirá salidas de ruido sin sentido. La regularización, al forzar que la distribución de las variables latentes se acerque a la normal estándar, rellena esos huecos.
\end{warningbox}

\subsection{Reparameterization Trick}

El entrenamiento del VAE enfrenta una dificultad técnica: la variable latente $z$ se obtiene mediante una operación de \textbf{muestreo}, y el muestreo no es diferenciable, por lo que no se puede retropropagar a través de la capa de muestreo.

\begin{figure}[H]
\centering
\includegraphics[width=0.85\textwidth]{slides-images/slide-35.jpg}
\caption{Reparameterization Trick: hace que el muestreo sea diferenciable\protect\footnotemark}
\end{figure}
\footnotetext{Slides página 35. Fuente: Kingma+ ICLR 2014.}

La solución es el \textbf{Reparameterization Trick} (truco de reparametrización). La idea central es separar la aleatoriedad de la variable latente $z$:

$$z = \mu + \sigma \odot \epsilon, \quad \epsilon \sim \mathcal{N}(0, 1)$$

\begin{itemize}
    \item $\mu$: el vector de medias que produce el codificador (nodo determinista diferenciable)
    \item $\sigma$: el vector de desviaciones estándar que produce el codificador (nodo determinista diferenciable)
    \item $\epsilon$: ruido aleatorio muestreado de la distribución normal estándar (constante aleatoria fija)
    \item $\odot$: multiplicación elemento a elemento
\end{itemize}

\begin{importantbox}{La esencia del Reparameterization Trick}
Al reescribir $z$ como una función determinista de $\mu$ y $\sigma$ más un ruido aleatorio externo $\epsilon$, el gradiente puede retropropagarse con normalidad a través de $\mu$ y $\sigma$. La aleatoriedad se "externaliza" como una constante muestreada de una distribución fija, y no bloquea el flujo del gradiente. Esto permite entrenar el VAE completo de extremo a extremo mediante retropropagación.
\end{importantbox}

\subsection{Perturbación de la variable latente e interpretabilidad}

Un VAE ya entrenado tiene una propiedad importante: las distintas dimensiones del espacio latente suelen corresponder a distintas \textbf{características interpretables}.

\begin{figure}[H]
\centering
\includegraphics[width=0.85\textwidth]{slides-images/slide-38.jpg}
\caption{Perturbación de la variable latente: fijar las demás variables y cambiar una sola para observar el cambio semántico\protect\footnotemark}
\end{figure}
\footnotetext{Slides página 38. Fuente: Kingma+ ICLR 2014.}

Por ejemplo, en un VAE entrenado con un conjunto de datos de rostros, si se fijan todas las variables latentes salvo una, cuyo valor se cambia lentamente, se puede observar un cambio suave en el rostro generado respecto a un atributo específico (como la orientación de la cabeza, la expresión, etc.). Esto muestra que el VAE logró aprender con éxito las características semánticas implícitas en los datos.

\subsection{$\beta$-VAE y el desenredo del espacio latente}

En el espacio latente de un VAE estándar, distintas dimensiones pueden codificar características entrelazadas entre sí. El $\beta$-VAE refuerza el desenredo agregando un hiperparámetro $\beta > 1$ delante del término de regularización:

$$\mathcal{L}(\theta, \phi; x, z, \beta) = \mathbb{E}_{q_\phi(z|x)}[\log p_\theta(x|z)] - \beta \cdot D_{KL}(q_\phi(z|x) \| p(z))$$

\begin{figure}[H]
\centering
\includegraphics[width=0.85\textwidth]{slides-images/slide-40.jpg}
\caption{El $\beta$-VAE logra un mejor desenredo del espacio latente\protect\footnotemark}
\end{figure}
\footnotetext{Slides página 40. Fuente: Higgins+ ICLR 2017.}

Cuando $\beta = 1$, se degrada al VAE estándar; cuando $\beta > 1$, se restringe con más fuerza el cuello de botella del espacio latente, lo que fomenta que cada variable latente codifique de manera independiente una característica semántica distinta (disentanglement), de modo que cambiar un atributo no afecte a los demás.

\subsection{Síntesis del VAE}

\begin{figure}[H]
\centering
\includegraphics[width=0.85\textwidth]{slides-images/slide-45.jpg}
\caption{Síntesis de los puntos clave del VAE\protect\footnotemark}
\end{figure}
\footnotetext{Slides página 45.}

Los cinco puntos clave del VAE:
\begin{enumerate}
    \item \textbf{Representación comprimida}: comprime la representación de alta dimensión del mundo en un espacio latente de baja dimensión que puede aprenderse
    \item \textbf{Aprendizaje no supervisado}: logra el aprendizaje sin etiquetas mediante la pérdida de reconstrucción
    \item \textbf{Reparameterization Trick}: permite entrenar de extremo a extremo redes que contienen una capa de muestreo
    \item \textbf{Variables latentes interpretables}: permite entender, mediante el análisis de perturbaciones, la característica semántica que corresponde a cada variable latente
    \item \textbf{Generación de muestras nuevas}: basta con muestrear del espacio latente y decodificar para generar datos nuevos
\end{enumerate}

\subsection{Resumen de la sección}

El VAE es una extensión probabilística del Autoencoder que, al introducir un espacio latente probabilístico y la regularización mediante KL Divergence, resuelve los problemas de discontinuidad e incompletitud del espacio latente del Autoencoder tradicional. El Reparameterization Trick resuelve con ingenio la dificultad de retropropagar a través de una capa de muestreo. El VAE es el primer modelo generativo profundo presentado en el curso, y pertenece a la familia de los Latent Variable Model.

